\documentclass[10pt,a4paper]{amsart}
\usepackage[margin=19mm]{geometry}
\usepackage{amsmath,amssymb,mathtools}
\usepackage[scr=rsfs,cal=euler]{mathalfa}
\usepackage{xfrac}
\usepackage[unicode,hidelinks,bookmarks=false]{hyperref}
\newcommand{\ie}{i.e.\ }
\newcommand{\cf}{cf.\ }
\newcommand{\resp}{respectively\ }
\newcommand{\Subsection}{\subsection}
\providecommand{\nobreakdash}{\nobreak\hbox{-}}
\setlength{\emergencystretch}{2em}
\multlinegap0pt
\usepackage{enumerate}
\usepackage[normalem]{ulem}
\usepackage{amssymb}
\usepackage{tikz}
\newtheorem{lem}{Lemma}
\def\dist{{\mathop\mathrm{\,dist\,}}}
\def\gz{{\gamma}}
\title{Sobolev trace inequalities on John domains and its applications}
\author{Weicong Su, Yi Ru-Ya Zhang}
\date{Source: arXiv:2406.06906v3 (CC BY 4.0)}
\begin{document}
\maketitle

\noindent\emph{Source and licence.} Sobolev trace inequalities on John domains and its applications. Version arXiv:2406.06906v3 (CC BY 4.0), distributed by arXiv under the Creative Commons Attribution 4.0 International licence, http://creativecommons.org/licenses/by/4.0/, as recorded in the arXiv licence metadata for this exact version and shown on the abstract page https://arxiv.org/abs/2406.06906v3. Full licence notice: \texttt{licenses/arxiv-2406.06906-CC-BY-4.0.txt}.\par\smallskip

\noindent\textit{Editorial scope.} Counted unit: the article's lem curve (source0.tex lines 622--624). The article's complete original proof of it is delivered with it (source0.tex lines 625--631, one \texttt{proof} environment of 7 lines). No mathematics is written, repaired or re-derived: the statement and the proof are the article's own lines, in the article's own notation, and no retained line is substituted or re-flowed.  Retained uncounted as the source-local context the proof needs: lines 307--309.  Nothing was omitted from any retained span, so the package carries no omission record.  No retained line contains a citation, so the wrapper carries no bibliography and no bibliography entry.  No retained line contains a figure environment, an image-inclusion command or drawing code, so no image is delivered and the package declares no asset dependency.  Editorial formatting only, in addition: an AMS \texttt{amsart} preamble that restates the article's own declarations and macros used by the retained lines, namely the environments \texttt{lem} on separate counters, together with the standard packages those lines need.  The retained environments print in this document as Lemma 1 in order of appearance, whereas the article prints the same items with the numbering of its own full document.  The complete native source of the article as distributed in the arXiv e-print for this exact version is bundled unchanged as \texttt{sources/160599/source0.tex}.
\begin{equation}\label{John curve}
\ell(\gamma[x,\,y])\le J\dist(y,\,\partial \Omega) \quad \text{ for any $y\in \gamma$}, 
\end{equation}

\begin{lem}\label{lem:boundary curve}
    Let $\Omega\subset\mathbb R^n$ be a John domain with John center $x_0\in \Omega$. Then for any $x\in \partial \Omega$, there exists (at least) one  curve $\gamma_x\subset \overline{\Omega}$ with $\gz_x\setminus\{x\}\subset \Omega$ from $x$ to $x_0$ for which \eqref{John curve} holds. 
\end{lem}

\begin{proof}
     Take a sequence of points $x_i\in \Omega$ approaching $x$,  and the corresponding John curve $\gz_{x_i}\subset \Omega$ joining $x$ to $x_0$. Then $\ell(\gz_{x_i})$ is uniformly bounded according to \eqref{John curve}. Now by parametrizing $\gamma_{x_i}$ via arc length and up to relabeling the sequence, Arzel\'a-Ascoli lemma yields the uniform convergence of  $\gz_i$ to some curve $\gz_x\subset \overline{\Omega}$ joining $x$ to $x_0$. 
     
Moreover, the uniform convergence also implies that $\gz_x$  satisfies \eqref{John curve}. This gives 
$$\gz_x\setminus\{x\}\subset \Omega.$$ 
This concludes the lemma.  
\end{proof}

\end{document}
